%----------------------------------------------------------------------------------------
% LIMIET VAN EEN RIJ. Alle opgaven zijn eigen werk.
%
% Opbouw: Start, leerstof (kern, uitgewerkt voorbeeld, korte oefening), en daarna op een
% nieuwe pagina de Opdrachten (door elkaar, stijgende moeilijkheid), Koppel en Test jezelf.
% Labels van dit hoofdstuk beginnen met h4-. \minigrafiek staat in hoofdstuk 1.
%----------------------------------------------------------------------------------------

\mychapter{4}{Limiet van een rij}

\lkr{Leerplan B+S'': LPD 45 (formele definitie, met kwantoren, toegepast op een concreet voorbeeld) en LPD 46 (limieten bepalen, convergentie, rekenregels, insluitstelling). De meetkundige reeks is een suggestie bij LPD 44 en staat hier omdat ze de limiet nodig heeft. Voorkennis: limieten van functies op oneindig. Rode draad: je kunt de limiet op drie manieren vinden, uit de grafiek, met rekenregels en uit een vergelijking, maar elke manier heeft haar voorwaarden.}

Een rij gaat altijd verder: na elke term komt er nog een. De vraag is dus niet wat de laatste term is, maar wat er op de lange duur gebeurt. Komen de termen steeds dichter bij een bepaald getal? Groeien ze onbegrensd? Of blijven ze schommelen? Je kent die vraag al van functies, waar je de limiet voor $x \to +\infty$ bepaalde. In dit hoofdstuk stel je ze voor rijen, en maak je eindelijk precies wat \textit{naderen} betekent. Daarmee kun je ook iets dat op het eerste gezicht onmogelijk lijkt: oneindig veel getallen optellen.

\begin{oefeningen}[Start]
\oef{h4-start-1}

\oef{h4-verken}
\end{oefeningen}

\lkr{Laat oefening \ref{oef:h4-verken} maken vóór de uitleg. Vraag c is de kern van de formele definitie: die komt straks in sectie 2 terug. Stel ook de vraag wat `naderen' eigenlijk betekent.}

%========================================================================================
\section{Convergent en divergent}

\subsection{De begrippen}

In de verkenning kwamen de termen steeds dichter bij een getal, en bleven ze er ook dicht bij. Zo'n getal heet de limiet van de rij. Niet elke rij heeft er een.

\begin{theorie}
Een rij is \emph{convergent} als de termen willekeurig dicht bij een reëel getal $L$ komen en daar blijven. $L$ is de \emph{limiet} van de rij: \ $\formule{\displaystyle\lim_{n\to+\infty} u_n = L}$. Een rij die niet convergeert, is \emph{divergent}.
\end{theorie}

Een divergente rij kan op twee manieren mislopen. Ze kan onbegrensd groeien (of dalen): dan schrijven we $\lim_{n\to+\infty} u_n = +\infty$ (of $-\infty$). Of ze kan blijven schommelen zonder ergens heen te gaan: dan bestaat de limiet niet.

\figuren{\minigrafiek{limiet 2}{-.5}{5.5}{2,4}{1/5,2/3.5,3/3,4/2.75,5/2.6,6/2.5}\qquad\qquad
\minigrafiek{limiet $+\infty$}{-.5}{10}{5,10}{1/.5,2/2,3/3,4/5,5/7,6/9.5}\qquad\qquad
\minigrafiek{geen limiet}{-3}{3}{-2,2}{1/-2,2/2,3/-2,4/2,5/-2,6/2}}

De limiet hoeft niet bereikt te worden, maar het mag wel: de rij $5,\ 5,\ 5,\ \ldots$ heeft limiet 5.

\voorbeeld{}{Bepaal de limiet. Is de rij convergent of divergent?\\[.5ex]
\textbf{a.}~$u_n = 4 - \tfrac1n$ \qquad \textbf{b.}~$u_n = 3n - n^2$ \qquad \textbf{c.}~$u_n = 1 + (-1)^n$\label{vb:h4-soort}}{%
\uitleg{Drie rijen die elk op een andere manier verlopen. Bij elke rij vragen we ons af wat er met de termen gebeurt als $n$ heel groot wordt.}
\begin{stappen}
\text{a. } \lim_{n\to+\infty} \left(4 - \dfrac1n\right) = 4 & De breuk $\tfrac1n$ wordt willekeurig klein, dus de termen komen willekeurig dicht bij 4: de rij is convergent. \\[1.5ex]
\text{b. } \lim_{n\to+\infty} \left(3n - n^2\right) = -\infty & Voor grote $n$ is $n^2$ veel groter dan $3n$, dus de term $-n^2$ beslist: de rij is divergent. \\[1ex]
\text{c. } 0,\ 2,\ 0,\ 2,\ \ldots & De termen springen heen en weer, en komen nooit dicht bij één getal. De limiet bestaat niet: de rij is divergent.
\end{stappen}}

\lkr{Veel leerlingen denken dat een limiet nooit bereikt wordt, en dat divergent hetzelfde is als `naar oneindig gaan'. De constante rij en de rij $1 + (-1)^n$ weerleggen allebei.}

\nujij{Bepaal de limiet van $u_n = \dfrac{(-1)^n}{n}$ en van $u_n = (-1)^n \cdot n$.}{Eerste rij: 0, convergent. Tweede rij: de limiet bestaat niet, divergent.}

\begin{oefeningen}
\oef{h4-k1}
\end{oefeningen}

\subsection{Rekenkundige en meetkundige rijen}

Bij rekenkundige en meetkundige rijen kun je de limiet meteen aflezen uit $v$ of $q$. De grafieken uit hoofdstuk 2 tonen het al: een rechte die stijgt of daalt, en een exponentiële kromme die groeit of naar 0 krimpt.

\begin{theorie}
\begin{stelling}{Limiet van een rekenkundige en een meetkundige rij}
{\renewcommand{\arraystretch}{1.35}%
\begin{tabular}{|>{\centering\arraybackslash}p{2.3cm}|>{\centering\arraybackslash}p{2.3cm}||>{\centering\arraybackslash}p{2.6cm}|>{\centering\arraybackslash}p{4.6cm}|}
\hline
\multicolumn{2}{|c||}{\textbf{rekenkundige rij}} & \multicolumn{2}{c|}{\textbf{meetkundige rij} ($u_1 \neq 0$)} \\ \hline
$v > 0$ & $\formule{+\infty}$ & $-1 < q < 1$ & $\formule{0}$ \\ \hline
$v < 0$ & $\formule{-\infty}$ & $q = 1$ & $\formule{u_1}$ \\ \hline
$v = 0$ & $\formule{u_1}$ & $q > 1$ & $\formule{+\infty}$ als $u_1 > 0$, $\formule{-\infty}$ als $u_1 < 0$ \\ \hline
 & & $q \leqslant -1$ & bestaat niet \\ \hline
\end{tabular}}
\end{stelling}

\begin{bewijs}
Een rekenkundige rij is $u_n = v\,n + (u_1 - v)$: de limiet van een eerstegraadsfunctie met richtingscoëfficiënt $v$. Een meetkundige rij is $u_n = u_1 \cdot q^{\,n-1}$. Bij $q > 0$ ken je de limiet van de exponentiële functie $q^x$: 0 als $q < 1$, en $+\infty$ als $q > 1$. Bij $-1 < q < 0$ is $|u_n| = |u_1| \cdot |q|^{\,n-1}$, en dat gaat naar 0, dus ook $u_n$. Bij $q \leqslant -1$ wisselen de termen van teken, terwijl $|u_n| \geqslant |u_1|$: ze kunnen dus niet dicht bij één getal blijven.
\end{bewijs}
\end{theorie}

\voorbeeld{}{Bepaal de limiet.\\[.5ex]
\textbf{a.}~$9,\ 7,\ 5,\ \ldots$ \qquad \textbf{b.}~$-3,\ {-6},\ {-12},\ \ldots$ \qquad \textbf{c.}~$8,\ {-4},\ 2,\ \ldots$ \qquad \textbf{d.}~$1,\ {-3},\ 9,\ \ldots$\label{vb:h4-rrmr}}{%
\uitleg{Eén rij per regel van de tabel. We bepalen telkens eerst de soort rij en $v$ of $q$, en zoeken dan de juiste regel.}
\begin{stappen}
\text{a. } v = -2 < 0 & Rekenkundig en dalend: de limiet is $-\infty$. \\
\text{b. } q = 2 > 1 \text{ en } u_1 < 0 & Meetkundig, en de termen worden steeds negatiever: de limiet is $-\infty$. \\
\text{c. } q = -\tfrac12 & Meetkundig met $-1 < q < 1$: de limiet is 0, ook al wisselt het teken. \\
\text{d. } q = -3 \leqslant -1 & Meetkundig, met tekens die wisselen en termen die groter worden: de limiet bestaat niet.
\end{stappen}}

\begin{oefeningen}
\oef{h4-k2}
\end{oefeningen}

%========================================================================================
\section{De formele definitie}

`Willekeurig dicht' klinkt duidelijk, maar wat betekent het precies? Hoe dicht is dicht genoeg? Wiskundigen hebben daar in de 19de eeuw een nauwkeurige definitie voor gevonden. Kies een band rond $L$, zo smal als je wilt, met breedte $\varepsilon$ (de Griekse letter epsilon) boven en onder $L$. Dan moeten vanaf een bepaald volgnummer $N$ alle termen in die band liggen. Kies je de band smaller, dan mag die $N$ groter worden, maar er moet er altijd een zijn.

\begin{theorie}
\begin{minipage}[t]{.57\linewidth}
\vspace{0pt}
$\formule{\displaystyle\lim_{n\to+\infty} u_n = L \ \Leftrightarrow}$

\vspace{.5mm}
$\formule{\forall\, \varepsilon > 0,\ \exists\, N \in \mathbb{N}: \ n > N \Rightarrow |u_n - L| < \varepsilon}$

\vspace{2.5mm}
$\formule{\displaystyle\lim_{n\to+\infty} u_n = +\infty \ \Leftrightarrow}$

\vspace{.5mm}
$\formule{\forall\, r > 0,\ \exists\, N \in \mathbb{N}: \ n > N \Rightarrow u_n > r}$

\vspace{2.5mm}
Bij $-\infty$ staat er op het einde $\formule{u_n < -r}$.
\end{minipage}\hfill
\begin{minipage}[t]{.4\linewidth}
\vspace{0pt}
\begin{tikzpicture}[x=.33cm,y=.72cm,>={Stealth[length=1.6mm]},font=\scriptsize]
  \fill[accent!18] (0,1.4) rectangle (13,2.6);
  \draw[lijngrijs,dashed] (0,2) -- (13,2);
  \draw[->,lijngrijs] (-.1,0) -- (13.5,0) node[right,text=black] {$n$};
  \draw[->,lijngrijs] (0,-.1) -- (0,5.4) node[above,text=black] {$u_n$};
  \draw[lijngrijs,dashed] (5,0) -- (5,2.6);
  \node[below=-1pt] at (5,0) {$N$};
  \node[left=-1pt] at (0,2) {$L$};
  \node[left=-1pt] at (0,2.75) {$L+\varepsilon$};
  \node[left=-1pt] at (0,1.25) {$L-\varepsilon$};
  \foreach \x/\y in {1/5,2/3.5,3/3,4/2.75,5/2.6,6/2.5,7/2.43,8/2.38,9/2.33,10/2.3,11/2.27,12/2.25}
    \fill[accent] (\x,\y) circle (1.5pt);
\end{tikzpicture}
\end{minipage}
\end{theorie}

Lees de definitie als een spel. Je tegenstander kiest een $\varepsilon$, hoe klein ook. Jij moet dan een $N$ vinden zodat alle termen na $u_N$ minder dan $\varepsilon$ van $L$ liggen. Lukt dat voor elke $\varepsilon$, dan is $L$ de limiet. De volgorde telt: $N$ mag van $\varepsilon$ afhangen. Bij een limiet $+\infty$ kiest je tegenstander geen band, maar een hoogte $r$, en moeten alle termen vanaf een bepaald volgnummer boven die hoogte liggen.

\voorbeeld{}{Bewijs met de definitie: $\displaystyle\lim_{n\to+\infty} \dfrac{2n+3}{n} = 2$.\label{vb:h4-bewijs}}{%
\uitleg{We moeten voor elke $\varepsilon$ een $N$ vinden. Daarvoor werken we de afstand tot de limiet uit, en zoeken we voor welke $n$ die kleiner is dan $\varepsilon$. Dat is kladwerk: het eigenlijke bewijs is de laatste stap.}
\begin{stappen}
\left|\dfrac{2n+3}{n} - 2\right| = \left|\dfrac{3}{n}\right| = \dfrac{3}{n} & Neem een willekeurige $\varepsilon > 0$. We werken de afstand tot 2 uit. \\[1.5ex]
\dfrac{3}{n} < \varepsilon \ \Leftrightarrow \ n > \dfrac{3}{\varepsilon} & We lossen op naar $n$. \\[1.5ex]
N \geqslant \dfrac{3}{\varepsilon},\ N \in \mathbb{N} & Kies zo'n $N$. Voor elke $n > N$ is dan $n > \tfrac{3}{\varepsilon}$, en dus $\left|u_n - 2\right| < \varepsilon$. Dat bewijst de limiet. \\[1.5ex]
\varepsilon = 0{,}01 \ \Rightarrow \ N = 300 & Bij een concrete $\varepsilon$ vind je een concrete $N$: vanaf $n = 301$ liggen alle termen in de band. Vergelijk met oefening \ref{oef:h4-verken}.
\end{stappen}}

\lkr{Laat het verschil zien tussen het kladwerk (oplossen naar $n$) en het bewijs zelf (kies $N$, en dan volgt het). Een concreet voorbeeld volstaat volgens het leerplan; de oefeningen op niveau 3 gaan verder.}

\voorbeeld{}{Bewijs met de definitie: $\displaystyle\lim_{n\to+\infty} \sqrt{n} = +\infty$.\label{vb:h4-bewijsinf}}{%
\uitleg{Nu zoeken we bij elke hoogte $r$ een $N$ waarna alle termen boven $r$ liggen.}
\begin{stappen}
\sqrt{n} > r \ \Leftrightarrow \ n > r^2 & Neem een willekeurige $r > 0$. Omdat beide leden positief zijn, mogen we kwadrateren. \\
N \geqslant r^2,\ N \in \mathbb{N} & Kies zo'n $N$. Voor elke $n > N$ is dan $n > r^2$, en dus $\sqrt{n} > r$. Dat bewijst de limiet.
\end{stappen}}

\begin{werkwijze}
Een limiet bewijzen met de definitie: neem een willekeurige $\varepsilon > 0$ (of $r > 0$), werk $|u_n - L| < \varepsilon$ (of $u_n > r$) uit tot een voorwaarde $n > \ldots$, en kies $N$ minstens zo groot. Schrijf daarna op dat voor elke $n > N$ de ongelijkheid geldt.
\end{werkwijze}

\nujij{Bewijs met de definitie: $\displaystyle\lim_{n\to+\infty} \dfrac{5}{n} = 0$.}{$\tfrac5n < \varepsilon$ geeft $n > \tfrac{5}{\varepsilon}$. Kies $N \geqslant \tfrac{5}{\varepsilon}$.}

\begin{oefeningen}
\oef{h4-k3}
\end{oefeningen}

%========================================================================================
\section{Limieten berekenen}

\subsection{Rekenregels}

Voor elke limiet een bewijs met $\varepsilon$ schrijven, is veel werk. Gelukkig hoeft dat niet: als de rij een voorschrift heeft dat ook voor alle reële $x$ zin heeft, mag je de limiet van die functie gebruiken. De functie neemt in de natuurlijke getallen precies de waarden van de rij aan.

\begin{theorie}
\begin{stelling}{Limiet van een rij en van een functie}
Is $u_n = f(n)$ en $\displaystyle\lim_{x\to+\infty} f(x) = L$, dan is $\formule{\displaystyle\lim_{n\to+\infty} u_n = L}$.
\end{stelling}

\begin{bewijs}
Neem $\varepsilon > 0$. Omdat $\lim_{x\to+\infty} f(x) = L$, is er een $X$ met $|f(x) - L| < \varepsilon$ voor elke $x > X$. Kies een natuurlijk getal $N \geqslant X$. Voor elke $n > N$ is dan ook $n > X$, dus $|u_n - L| = |f(n) - L| < \varepsilon$.
\end{bewijs}
\end{theorie}

Zo mag je alle rekenregels voor limieten van functies gebruiken: de limiet van een som, een verschil, een product en een quotiënt. De rijen $\tfrac1n$ en $\tfrac{1}{n^2}$ hebben bijvoorbeeld limiet 0.

Het omgekeerde geldt niet. De rij $u_n = \sin(n\pi)$ is $0,\ 0,\ 0,\ \ldots$ met limiet 0, maar $\lim_{x\to+\infty} \sin(x\pi)$ bestaat niet: tussen de natuurlijke getallen blijft de sinus schommelen. Let ook op met $\infty - \infty$: dat is niet 0, maar een onbepaaldheid. Reken eerst om.

\voorbeeld{}{Bereken $\displaystyle\lim_{n\to+\infty} \dfrac{3n^2 - n}{n^2 + 2}$.\label{vb:h4-reken}}{%
\uitleg{Teller en noemer gaan allebei naar $+\infty$. Daarom delen we ze door de hoogste macht van $n$ die in de noemer staat: dan blijven er termen over die naar 0 gaan.}
\begin{stappen}
\dfrac{3n^2 - n}{n^2 + 2} = \dfrac{3 - \tfrac1n}{1 + \tfrac{2}{n^2}} & We delen teller en noemer door $n^2$. \\[3ex]
\lim_{n\to+\infty} \dfrac{3 - \tfrac1n}{1 + \tfrac{2}{n^2}} = \dfrac{3 - 0}{1 + 0} = 3 & De breuken $\tfrac1n$ en $\tfrac{2}{n^2}$ gaan naar 0.
\end{stappen}}

\voorbeeld{}{Bereken $\displaystyle\lim_{n\to+\infty} \left(\sqrt{n+1} - \sqrt{n}\right)$.\label{vb:h4-reken2}}{%
\uitleg{Beide wortels gaan naar $+\infty$: dit is de onbepaaldheid $\infty - \infty$. Met de toegevoegde uitdrukking maken we er een breuk van, zodat de wortels niet meer van elkaar afgetrokken worden.}
\begin{stappen}
\sqrt{n+1} - \sqrt{n} = \dfrac{(n+1) - n}{\sqrt{n+1} + \sqrt{n}} & We vermenigvuldigen teller en noemer met $\sqrt{n+1} + \sqrt{n}$. In de teller staat dan een merkwaardig product. \\[3ex]
\lim_{n\to+\infty} \dfrac{1}{\sqrt{n+1} + \sqrt{n}} = 0 & De teller is 1, en de noemer wordt onbegrensd groot.
\end{stappen}}

\nujij{Bereken $\displaystyle\lim_{n\to+\infty} \dfrac{2n - n^2}{4n^2 + 1}$.}{$-\tfrac14$: deel teller en noemer door $n^2$.}

\begin{oefeningen}
\oef{h4-k4}
\end{oefeningen}

\subsection{De insluitstelling}

Sommige rijen bevatten een factor die blijft schommelen, zoals $\sin n$, $\cos n$ of $(-1)^n$. Daarvan bestaat de limiet niet, dus de rekenregels helpen niet. Zo'n factor is wel begrensd. Kun je de rij dan insluiten tussen twee rijen met dezelfde limiet, dan heeft ze geen andere keuze.

\begin{theorie}
\begin{stelling}{Insluitstelling}
Als $t_n \leqslant u_n \leqslant v_n$ vanaf een bepaald volgnummer, en $\displaystyle\lim_{n\to+\infty} t_n = \lim_{n\to+\infty} v_n = L$, dan is ook $\formule{\displaystyle\lim_{n\to+\infty} u_n = L}$.
\end{stelling}

\begin{bewijs}
Neem $\varepsilon > 0$. Vanaf een bepaald volgnummer ligt $t_n$ boven $L - \varepsilon$, vanaf een ander volgnummer ligt $v_n$ onder $L + \varepsilon$, en vanaf een derde geldt $t_n \leqslant u_n \leqslant v_n$. Neem voor $N$ het grootste van de drie. Voor elke $n > N$ is dan
\bewijsregel{L - \varepsilon < t_n \leqslant u_n \leqslant v_n < L + \varepsilon , \quad\text{dus}\quad |u_n - L| < \varepsilon .}
\end{bewijs}
\end{theorie}

Informeel gezegd: zit een rij geklemd tussen twee rijen die naar hetzelfde getal gaan, dan moet ze mee.

\voorbeeld{}{Bereken $\displaystyle\lim_{n\to+\infty} \dfrac{\cos n}{n}$.\label{vb:h4-insluit}}{%
\naast{\uitleg{De limiet van $\cos n$ bestaat niet, maar $\cos n$ ligt altijd tussen $-1$ en 1. We sluiten de rij in tussen twee eenvoudige rijen.}
\begin{stappenvrij}{4cm}{4.6cm}
-1 \leqslant \cos n \leqslant 1 & De cosinus is begrensd. \\[1ex]
-\dfrac1n \leqslant \dfrac{\cos n}{n} \leqslant \dfrac1n & We delen door $n > 0$: de tekens blijven. \\[1.5ex]
\lim_{n\to+\infty} \dfrac{\cos n}{n} = 0 & Beide grenzen hebben limiet 0, dus ook de ingesloten rij.
\end{stappenvrij}}{\begin{tikzpicture}[x=.5cm,y=1.05cm,font=\scriptsize]
  \draw[black!14,very thin,xstep=1,ystep=.5] (0,-1) grid (8.4,1);
  \draw[->,grijs] (0,0) -- (8.8,0) node[right,text=black] {$n$};
  \draw[->,grijs] (0,-1.1) -- (0,1.2) node[above,text=black] {$u_n$};
  \foreach \t in {2,4,6,8}{\node[below,inner sep=1.5pt] at (\t,-1) {$\t$};}
  \foreach \t in {-1,1}{\node[left,inner sep=1.5pt] at (0,\t) {$\t$};}
  \draw[grijs,dashed,domain=.9:8.4,samples=40] plot (\x,{1/\x});
  \draw[grijs,dashed,domain=.9:8.4,samples=40] plot (\x,{-1/\x});
  \node[right=-1pt] at (1.25,.82) {$\tfrac1n$};
  \node[right=-1pt] at (1.25,-.82) {$-\tfrac1n$};
  \foreach \x/\y in {1/.54,2/-.208,3/-.33,4/-.163,5/.057,6/.16,7/.108,8/-.018}
    \fill[accent] (\x,\y) circle (1.6pt);
\end{tikzpicture}}}

\begin{oefeningen}
\oef{h4-k5}
\end{oefeningen}

%========================================================================================
\section{Stijgende en dalende rijen}

Soms kun je een limiet niet uitrekenen, maar wil je toch weten of ze bestaat. Dan helpt het verloop van de rij: gaan de termen alleen maar omhoog, en is er een plafond dat ze nooit doorbreken?

\begin{theorie}
Een rij is \emph{stijgend} als $\formule{u_{n+1} > u_n}$ voor elke $n$, en \emph{dalend} als $\formule{u_{n+1} < u_n}$ voor elke $n$.

\vspace{1mm}
Een rij is \emph{naar boven begrensd} als er een $M$ is met $\formule{u_n \leqslant M}$ voor elke $n$, en \emph{naar onder begrensd} als er een $m$ is met $\formule{u_n \geqslant m}$ voor elke $n$. Is ze allebei, dan is ze \emph{begrensd}.

\begin{stelling}{Stijgend en begrensd}
Een stijgende, naar boven begrensde rij is convergent. Een dalende, naar onder begrensde rij ook.
\end{stelling}
\end{theorie}

Stel je de grafiek voor: de punten gaan alleen maar omhoog, maar er is een plafond dat ze niet kunnen doorbreken. Dan moeten ze zich wel ophopen onder een bepaalde hoogte, en die hoogte is de limiet. Het bewijs steunt op een eigenschap van de reële getallen die je in het secundair niet ziet: we gebruiken de stelling dus zonder bewijs. Eén voorwaarde volstaat niet: $u_n = n$ is stijgend en $u_n = (-1)^n$ is begrensd, maar beide zijn divergent.

\voorbeeld{}{Toon aan dat $u_n = \dfrac{n}{n+2}$ stijgend en naar boven begrensd is. Wat is de limiet?\label{vb:h4-monotoon}}{%
\uitleg{Voor de bovengrens herschrijven we de breuk, zodat we zien dat ze kleiner is dan 1. Om te tonen dat de rij stijgt, berekenen we het verschil van twee opeenvolgende termen, en tonen we dat het positief is.}
\begin{stappen}
u_n = \dfrac{(n+2) - 2}{n+2} = 1 - \dfrac{2}{n+2} & We trekken 2 af in de teller, en tellen het er weer bij. Dus $u_n < 1$ voor elke $n$: 1 is een bovengrens. \\[2ex]
u_{n+1} - u_n = \dfrac{2}{n+2} - \dfrac{2}{n+3} = \dfrac{2}{(n+2)(n+3)} & Het verschil van twee opeenvolgende termen, op gelijke noemer gezet. \\[2ex]
u_{n+1} - u_n > 0 & Teller en noemer zijn positief voor elke $n$: de rij stijgt. Volgens de stelling is ze dus convergent. \\[1ex]
\lim_{n\to+\infty} u_n = 1 & De stelling zegt dat er een limiet is. De rekenregels geven de waarde.
\end{stappen}
\naast{De figuur toont het: de punten stijgen, maar blijven onder het plafond 1, en kruipen er steeds dichter tegenaan.}{\begin{tikzpicture}[x=.4cm,y=2cm,font=\scriptsize]
  \draw[black!14,very thin,xstep=1,ystep=.25] (0,0) grid (10.4,1.1);
  \draw[->,grijs] (0,0) -- (10.8,0) node[right,text=black] {$n$};
  \draw[->,grijs] (0,0) -- (0,1.2) node[above,text=black] {$u_n$};
  \foreach \t in {2,4,...,10}{\node[below,inner sep=1.5pt] at (\t,0) {$\t$};}
  \node[left,inner sep=1.5pt] at (0,.5) {$0{,}5$};
  \node[left,inner sep=1.5pt] at (0,1) {$1$};
  \draw[grijs,dashed] (0,1) -- (10.4,1);
  \foreach \x in {1,...,10}{\fill[accent] (\x,{\x/(\x+2)}) circle (1.6pt);}
\end{tikzpicture}}}

\begin{oefeningen}
\oef{h4-k6}
\end{oefeningen}

%========================================================================================
\section{Discrete dynamische processen}

In opdracht \ref{oef:h1-o14} van hoofdstuk 1 verdween er elke nacht 20\,\% van het chloor in een zwembad, en kwam er elke ochtend 50 g bij. Zo'n rij met $u_n = a\,u_{n-1} + b$ heet een \textit{discreet dynamisch proces}. Discreet: je kijkt met vaste tussenpozen. Dynamisch: er verandert iets. Op de lange duur komt zo'n proces vaak tot rust in een \textit{evenwicht}: er verdwijnt dan evenveel als er bijkomt.

\begin{theorie}
\begin{stelling}{Evenwicht van een dynamisch proces}
Is $u_n = a\,u_{n-1} + b$ met $|a| < 1$, dan is de rij convergent, en haar limiet $L$ voldoet aan

\vspace{1mm}
\qquad $\formule{L = a L + b}$ \quad dus \quad $\formule{L = \dfrac{b}{1 - a}}$
\end{stelling}

\begin{bewijs}
Neem $L = \tfrac{b}{1 - a}$, zodat $L = aL + b$, en bekijk de hulprij $v_n = u_n - L$. Trek de twee gelijkheden af:
\bewijsregel{v_n = u_n - L = (a\,u_{n-1} + b) - (aL + b) = a\,(u_{n-1} - L) = a\,v_{n-1} .}
De hulprij is meetkundig met quotiënt $a$, en omdat $|a| < 1$, gaat ze naar 0. Dan gaat $u_n = L + v_n$ naar $L$.
\end{bewijs}
\end{theorie}

\voorbeeld{}{Het zwembad: $u_0 = 400$ en $u_n = 0{,}8\,u_{n-1} + 50$. Bepaal het evenwicht $L$ en een expliciet voorschrift.\label{vb:h4-dyn}}{%
\naast{\uitleg{Omdat $|0{,}8| < 1$, bestaat de limiet. We vinden ze uit de vergelijking, en het expliciet voorschrift met de hulprij uit het bewijs.}
\begin{stappenvrij}{4.6cm}{4.2cm}
L = 0{,}8\,L + 50 \ \Rightarrow \ L = 250 & $0{,}2\,L = 50$: het evenwicht is 250 g chloor. \\
v_n = u_n - 250 & De hulprij is meetkundig met $q = 0{,}8$ en $v_0 = 400 - 250 = 150$. \\
u_n = 250 + 150 \cdot 0{,}8^{\,n} & Dus $v_n = 150 \cdot 0{,}8^{\,n}$. Controle: $u_1 = 250 + 120 = 370$. \\
\lim_{n\to+\infty} u_n = 250 & Want $0{,}8^{\,n} \to 0$.
\end{stappenvrij}}{\begin{tikzpicture}[x=.38cm,y=.0095cm,font=\scriptsize]
  \draw[black!14,very thin,xstep=1,ystep=50] (0,0) grid (10.4,420);
  \draw[->,grijs] (0,0) -- (10.8,0) node[right,text=black] {$n$};
  \draw[->,grijs] (0,0) -- (0,450) node[above,text=black] {$u_n$};
  \foreach \t in {2,4,...,10}{\node[below,inner sep=1.5pt] at (\t,0) {$\t$};}
  \foreach \t in {100,200,300,400}{\node[left,inner sep=1.5pt] at (0,\t) {$\t$};}
  \draw[grijs,dashed] (0,250) -- (10.4,250) node[right,text=black] {250};
  \foreach \x/\y in {0/400,1/370,2/346,3/326.8,4/311.4,5/299.2,6/289.3,7/281.5,8/275.2,9/270.1,10/266.1}
    \fill[accent] (\x,\y) circle (1.6pt);
\end{tikzpicture}}}

De vergelijking $L = aL + b$ veronderstelt dat de limiet bestaat. Is $|a| \geqslant 1$, dan levert ze vaak toch een getal op, terwijl de rij geen limiet heeft. Kijk dus altijd eerst naar $a$.

\lkr{De hulprij $v_n$ sluit aan bij opdracht \ref{oef:h1-o21} van hoofdstuk 1. In de opdrachten staat ook een geval met $|a| \geqslant 1$.}

\nujij{$u_0 = 90$ en $u_n = 0{,}5\,u_{n-1} + 10$. Bepaal het evenwicht en een expliciet voorschrift.}{$L = 20$ en $u_n = 20 + 70 \cdot 0{,}5^{\,n}$.}

\begin{oefeningen}
\oef{h4-k7}
\end{oefeningen}

%========================================================================================
\section{De meetkundige reeks}

Een slak kruipt een meter, dan een halve meter, dan een kwart meter, en zo verder. Komt ze ooit voorbij twee meter? De afgelegde weg na $n$ dagen is de som $s_n = 1 + \tfrac12 + \cdots + \tfrac{1}{2^{n-1}} = 2 - \tfrac{1}{2^{n-1}}$, en die komt willekeurig dicht bij 2. Zo kunnen oneindig veel getallen samen toch een eindige som hebben.

\begin{theorie}
Een \emph{reeks} is een oneindige som. Haar waarde is de limiet van de sommen $s_n$, als die bestaat.

\begin{stelling}{Meetkundige reeks}
Voor een meetkundige rij met $|q| < 1$ is

\vspace{1mm}
\qquad $\formule{u_1 + u_1 q + u_1 q^2 + \cdots = \dfrac{u_1}{1 - q}}$

\vspace{1mm}
Bij $|q| \geqslant 1$ en $u_1 \neq 0$ heeft de reeks geen eindige waarde: ze \emph{divergeert}.
\end{stelling}

\begin{bewijs}
Volgens de somformule is $s_n = u_1 \cdot \dfrac{q^{\,n} - 1}{q - 1} = u_1 \cdot \dfrac{1 - q^{\,n}}{1 - q}$. Bij $|q| < 1$ gaat $q^{\,n}$ naar 0, zodat
\bewijsregel{\lim_{n\to+\infty} s_n = u_1 \cdot \dfrac{1 - 0}{1 - q} = \dfrac{u_1}{1 - q} .}
Bij $|q| \geqslant 1$ gaan de termen zelf niet naar 0, en dan kan de som niet tot rust komen.
\end{bewijs}
\end{theorie}

\voorbeeld{}{Bereken $8 + 2 + \tfrac12 + \cdots$\label{vb:h4-reeks}}{%
\uitleg{De termen worden telkens vier keer kleiner: dit is een meetkundige reeks. De formule mag alleen als $|q| < 1$.}
\begin{stappen}
q = \dfrac28 = \dfrac14 & De reeks is meetkundig, en $|q| < 1$. \\[1.5ex]
\dfrac{8}{1 - \tfrac14} = \dfrac{8}{\tfrac34} = \dfrac{32}{3} & De som van alle termen is $\tfrac{32}{3} \approx 10{,}67$.
\end{stappen}}

\voorbeeld{}{Schrijf $0{,}272727\ldots$ als een breuk.\label{vb:h4-reeks2}}{%
\uitleg{Een repeterend decimaal getal is een oneindige som: het blok 27 komt telkens terug, honderd keer kleiner.}
\begin{stappen}
0{,}27 + 0{,}0027 + 0{,}000027 + \cdots & We schrijven het getal als een reeks, met $u_1 = 0{,}27$ en $q = 0{,}01$. \\[1ex]
\dfrac{0{,}27}{1 - 0{,}01} = \dfrac{27}{99} = \dfrac{3}{11} & De formule voor de meetkundige reeks. Controle met het rekentoestel: $3 : 11 = 0{,}2727\ldots$
\end{stappen}}

\nujij{Bereken $9 - 3 + 1 - \cdots$ en schrijf $0{,}1515\ldots$ als een breuk.}{$\tfrac{9}{1 + 1/3} = \tfrac{27}{4}$ en $\tfrac{15}{99} = \tfrac{5}{33}$.}

\lkr{De slak kwam al voor in een opdracht van hoofdstuk 3. De reeks is een suggestie in het leerplan, geen doel.}

\begin{oefeningen}
\oef{h4-k8}
\end{oefeningen}

%----------------------------------------------------------------------------------------
% OPDRACHTEN: door elkaar. Niveau 1, dan 2, dan 3.
%----------------------------------------------------------------------------------------
\opdrachtenkop

\lkr{De opdrachten staan bewust door elkaar, met stijgende moeilijkheid. In deze versie staat bij elke opdracht het onderwerp, zodat je kunt kiezen. Vraag bij elke limiet ook: convergent of divergent?}

\begin{oefeningen}
\oef{h4-soort}

\oef{h4-reeks}

\oef{h4-ratio}

\oef{h4-rrmr}

\oef{h4-o1}

\oef{h4-monotoon}

\oef{h4-evenwicht}

\oef{h4-epsgetal}

\oef{h4-bewijs}

\oef{h4-o2}

\oef{h4-o3}

\oef{h4-insluit}

\oef{h4-o4}

\oef{h4-o5}

\oef{h4-o6}

\oef{h4-o7}

\oef{h4-o8}

\oef{h4-o9}

\oef{h4-o10}

\oef{h4-o11}

\oef{h4-o12}

\oef{h4-o13}

\oef{h4-o14}

\oef{h4-o15}

\oef{h4-o16}

\oef{h4-o17}
\end{oefeningen}

%----------------------------------------------------------------------------------------
% KOPPEL: verbanden met andere leerstof
%----------------------------------------------------------------------------------------
\opdrachtsectie{Koppel}

\lkr{Exponentiële functies, complexe getallen en matrices zijn gezien in het 5de jaar. Kijk na of het getal $e$ al ingevoerd is. Logica en kansrekenen staan ook in het 6de jaar: kijk na of ze al aan bod kwamen, of laat die oefeningen weg.}

\begin{oefeningen}
\oef{h4-koppel-1}

\oef{h4-koppel-2}

\oef{h4-koppel-3}

\oef{h4-koppel-4}

\oef{h4-koppel-5}
\end{oefeningen}

%----------------------------------------------------------------------------------------
\opdrachtsectie{Test jezelf}

\begin{oefeningen}
\oef{h4-t1}

\oef{h4-t2}

\oef{h4-t3}

\oef{h4-t4}

\oef{h4-t5}

\oef{h4-t6}

\oef{h4-t7}

\oef{h4-t8}
\end{oefeningen}

\herhaalschema{%
  \fv{oef:h4-t1}{vb:h4-soort}{oef:h4-soort}%
  \fv{oef:h4-t2}{vb:h4-rrmr}{oef:h4-rrmr}%
  \fv{oef:h4-t3}{vb:h4-bewijs}{oef:h4-bewijs}%
  \fv{oef:h4-t4}{vb:h4-reken}{oef:h4-ratio}%
  \fv{oef:h4-t5}{vb:h4-insluit}{oef:h4-insluit}%
  \fv{oef:h4-t6}{vb:h4-monotoon}{oef:h4-monotoon}%
  \fv{oef:h4-t7}{vb:h4-dyn}{oef:h4-evenwicht}%
  \fv{oef:h4-t8}{vb:h4-reeks}{oef:h4-reeks}}

\ifzelfstudie
\section*{Maak je eigen samenvatting}
Sluit dit boek en schrijf in je schrift: wat convergent en divergent betekent, de formele definitie, wanneer een rekenkundige en een meetkundige rij convergeren, hoe je een limiet berekent, de insluitstelling, de eigenschap van stijgende en begrensde rijen, hoe je een evenwicht vindt, en de som van een meetkundige reeks. Vergelijk daarna met de kaders.
\fi
